% ========================================================
% BÀI TẬP TOÁN 5 - KẾT NỐI TRI THỨC VỚI CUỘC SỐNG
% Bài 19: Phép cộng số thập phân
% Định dạng: BTTeX Part III - Trắc nghiệm điền số ngắn
% ========================================================

% Câu 1
\begin{ex}
  Tính kết quả: $14{,}5 + 25{,}5$.
  \shortans{40}
  \loigiai{
    $14{,}5 + 25{,}5 = 40$.
  }
\end{ex}

% Câu 2
\begin{ex}
  Tính kết quả: $38{,}2 + 61{,}8$.
  \shortans{100}
  \loigiai{
    $38{,}2 + 61{,}8 = 100$.
  }
\end{ex}

% Câu 3
\begin{ex}
  Tính tổng $23{,}45 + 16{,}55$.
  \shortans{40}
  \loigiai{
    $23{,}45 + 16{,}55 = 40$.
  }
\end{ex}

% Câu 4
\begin{ex}
  Tính tổng $125{,}3 + 74{,}7$.
  \shortans{200}
  \loigiai{
    $125{,}3 + 74{,}7 = 200$.
  }
\end{ex}

% Câu 5
\begin{ex}
  Tính nhanh: $15{,}8 + 24{,}2 + 10$.
  \shortans{50}
  \loigiai{
    $(15{,}8 + 24{,}2) + 10 = 40 + 10 = 50$.
  }
\end{ex}

% Câu 6
\begin{ex}
  Tính nhanh bằng cách hợp lí: $12{,}5 + 34{,}7 + 7{,}5 + 5{,}3$.
  \shortans{60}
  \loigiai{
    $(12{,}5 + 7{,}5) + (34{,}7 + 5{,}3) = 20 + 40 = 60$.
  }
\end{ex}

% Câu 7
\begin{ex}
  Tính nhanh: $4{,}85 + 7{,}2 + 5{,}15 + 2{,}8$.
  \shortans{20}
  \loigiai{
    $(4{,}85 + 5{,}15) + (7{,}2 + 2{,}8) = 10 + 10 = 20$.
  }
\end{ex}

% Câu 8
\begin{ex}
  Tìm số tự nhiên $x$ biết: $x - 14{,}5 = 25{,}5$.
  \shortans{40}
  \loigiai{
    $x = 25{,}5 + 14{,}5 = 40$.
  }
\end{ex}

% Câu 9
\begin{ex}
  Tìm $x$ biết: $x - 38{,}25 = 11{,}75$.
  \shortans{50}
  \loigiai{
    $x = 11{,}75 + 38{,}25 = 50$.
  }
\end{ex}

% Câu 10
\begin{ex}
  Một người thợ dệt ngày thứ nhất dệt được $28{,}4\text{ m}$ vải, ngày thứ hai dệt được $31{,}6\text{ m}$ vải. Cả hai ngày người đó dệt được bao nhiêu mét vải?
  \shortans{60}
  \loigiai{
    Cả hai ngày dệt được: $28{,}4 + 31{,}6 = 60\text{ m}$.
  }
\end{ex}

% Câu 11
\begin{ex}
  Một bao gạo cân nặng $45{,}5\text{ kg}$, một bao ngô cân nặng $34{,}5\text{ kg}$. Cả hai bao cân nặng bao nhiêu ki-lô-gam?
  \shortans{80}
  \loigiai{
    Cả hai bao nặng: $45{,}5 + 34{,}5 = 80\text{ kg}$.
  }
\end{ex}

% Câu 12
\begin{ex}
  Một hình tam giác có độ dài ba cạnh lần lượt là $4{,}5\text{ cm}; 5{,}5\text{ cm}$ và $6\text{ cm}$. Chu vi hình tam giác đó là bao nhiêu xăng-ti-mét?
  \shortans{16}
  \loigiai{
    Chu vi là: $4{,}5 + 5{,}5 + 6 = 16\text{ cm}$.
  }
\end{ex}

% Câu 13
\begin{ex}
  Một sợi dây dài $12{,}5\text{ m}$, người ta nối thêm một đoạn dây dài $7{,}5\text{ m}$. Chiều dài sợi dây sau khi nối là bao nhiêu mét?
  \shortans{20}
  \loigiai{
    Chiều dài mới là $12{,}5 + 7{,}5 = 20\text{ m}$.
  }
\end{ex}

% Câu 14
\begin{ex}
  Tìm số tự nhiên $a$ biết: $a = 15{,}7 + 24{,}3$.
  \shortans{40}
  \loigiai{
    $a = 15{,}7 + 24{,}3 = 40$.
  }
\end{ex}

% Câu 15
\begin{ex}
  Một thùng dầu chứa $18{,}25$ lít dầu, người ta đổ thêm vào thùng $11{,}75$ lít dầu. Hỏi trong thùng có tất cả bao nhiêu lít dầu?
  \shortans{30}
  \loigiai{
    Trong thùng có: $18{,}25 + 11{,}75 = 30$ lít.
  }
\end{ex}

% Câu 16
\begin{ex}
  Tính giá trị biểu thức: $3{,}25 + 4{,}75 + 12$.
  \shortans{20}
  \loigiai{
    $(3{,}25 + 4{,}75) + 12 = 8 + 12 = 20$.
  }
\end{ex}

% Câu 17
\begin{ex}
  Tổng của ba số là $45$. Biết số thứ nhất là $12{,}4$, số thứ hai là $17{,}6$. Số thứ ba là bao nhiêu?
  \shortans{15}
  \loigiai{
    Tổng hai số đầu là $12{,}4 + 17{,}6 = 30$. Số thứ ba là $45 - 30 = 15$.
  }
\end{ex}

% Câu 18
\begin{ex}
  Tính: $0{,}5 + 1{,}5 + 2{,}5 + 3{,}5 + 4{,}5 + 5{,}5$.
  \shortans{18}
  \loigiai{
    $(0{,}5 + 5{,}5) + (1{,}5 + 4{,}5) + (2{,}5 + 3{,}5) = 6 + 6 + 6 = 18$.
  }
\end{ex}

% Câu 19
\begin{ex}
  Một cửa hàng buổi sáng bán được $120{,}5\text{ kg}$ gạo, buổi chiều bán được $129{,}5\text{ kg}$ gạo. Cả ngày cửa hàng bán được bao nhiêu ki-lô-gam gạo?
  \shortans{250}
  \loigiai{
    $120{,}5 + 129{,}5 = 250\text{ kg}$.
  }
\end{ex}

% Câu 20
\begin{ex}
  Một mảnh vườn hình chữ nhật có chiều rộng $14{,}5\text{ m}$, chiều dài hơn chiều rộng $5\text{ m}$. Nửa chu vi mảnh vườn là bao nhiêu mét?
  \shortans{34}
  \loigiai{
    Chiều dài là $14{,}5 + 5 = 19{,}5\text{ m}$. Nửa chu vi là $14{,}5 + 19{,}5 = 34\text{ m}$.
  }
\end{ex}

% Câu 21
\begin{ex}
  Tìm chữ số $x$ để tổng $12{,}x + 15{,}4$ có chữ số tận cùng của phần thập phân là 0.
  \shortans{6}
  \loigiai{
    Để $x + 4 = 10$ thì $x = 6$.
  }
\end{ex}

% Câu 22
\begin{ex}
  Tính tổng của số thập phân lớn nhất có một chữ số ở phần thập phân bé hơn 10 với số thập phân bé nhất có một chữ số ở phần thập phân lớn hơn 0: $9{,}9 + 0{,}1$.
  \shortans{10}
  \loigiai{
    $9{,}9 + 0{,}1 = 10$.
  }
\end{ex}

% Câu 23
\begin{ex}
  Một can nước ngọt có khối lượng vỏ là $0{,}5\text{ kg}$, lượng nước ngọt bên trong nặng $9{,}5\text{ kg}$. Tổng khối lượng cả can là bao nhiêu ki-lô-gam?
  \shortans{10}
  \loigiai{
    $0{,}5 + 9{,}5 = 10\text{ kg}$.
  }
\end{ex}

% Câu 24
\begin{ex}
  Tính nhanh giá trị biểu thức: $2{,}4 + 7{,}6 + 3{,}8 + 6{,}2 + 5$.
  \shortans{25}
  \loigiai{
    $(2{,}4 + 7{,}6) + (3{,}8 + 6{,}2) + 5 = 10 + 10 + 5 = 25$.
  }
\end{ex}

% Câu 25
\begin{ex}
  Khi thực hiện phép cộng một số tự nhiên với một số thập phân có một chữ số ở phần thập phân, một bạn quên dấu phẩy nên được kết quả là $183$. Biết số thập phân có chữ số hàng phần mười là $3$. Số tự nhiên đó là bao nhiêu?
  \shortans{150}
  \loigiai{
    Số thập phân là $3{,}3$ quên dấu phẩy thành 33. Số tự nhiên là $183 - 33 = 150$.
  }
\end{ex}
